\chapter{Cost Curves and the Firm's Supply}
\label{chap:cost-curves}

\section{Learning objectives}

After completing this chapter, you should be able to:
\begin{itemize}
  \item distinguish fixed, variable, total, average, and marginal costs;
  \item explain the shapes of standard short-run cost curves;
  \item apply the shutdown rule; and
  \item derive the short-run supply curve of a competitive firm.
\end{itemize}

\section{Short-run costs}

In the short run, some inputs are fixed. \keyterm{Fixed cost} ($FC$) does not vary with output: rent on a lease is a common example. \keyterm{Variable cost} ($VC(q)$) changes with output, such as the cost of ingredients or hourly labour. Total cost is
\[
TC(q)=FC+VC(q).
\]

Fixed cost is paid even when output is zero; variable cost is zero when no variable inputs are used. This distinction is about whether a cost can be avoided by producing nothing in the relevant period, not about whether a payment is large or small.

\section{Average and marginal costs}

Average fixed cost, average variable cost, and average total cost are
\[
AFC=\frac{FC}{q}, \qquad AVC=\frac{VC(q)}{q}, \qquad ATC=\frac{TC(q)}{q}=AFC+AVC.
\]
Marginal cost is the additional cost of producing an additional unit:
\[
MC=\frac{dTC}{dq}=\frac{dVC}{dq}.
\]
Fixed cost does not affect marginal cost because it does not change with output.

Average fixed cost always falls as output rises, because a fixed amount is spread over more units. Average variable cost and average total cost are often U-shaped. At low output, specialisation can reduce average cost; at higher output, diminishing marginal product raises the variable input needed for extra output.

\section{The marginal--average relationship}

Marginal cost crosses both AVC and ATC at their minimum points. The intuition is general: a marginal value below an average pulls the average down, while a marginal value above an average pushes it up. For example, if the cost of the next unit is below current average cost, producing it lowers the average.

\section{The competitive firm's output choice}

A competitive firm is a \keyterm{price taker}: it can sell as much as it wishes at the market price but cannot influence that price. Total revenue is $TR=pq$, so
\[
MR=\frac{dTR}{dq}=p.
\]
The firm chooses output where $p=MC$, provided it is worthwhile to operate. The condition identifies the output level at which the additional revenue from another unit equals its additional cost.

\section{The shutdown decision}

In the short run, fixed costs are unavoidable. A firm should produce if revenue can cover variable cost, or equivalently if price is at least minimum average variable cost:
\[
p\geq \min AVC.
\]
If $p<AVC$, each unit produced adds more variable cost than revenue, so shutting down minimises the loss. If $AVC\leq p<ATC$, the firm produces but makes a loss; it covers variable cost and contributes something toward fixed cost. In the long run, all costs are avoidable, so a firm remains in the market only if price covers average total cost.

\section{Supply}

The competitive firm's short-run supply curve is the portion of its marginal-cost curve above the minimum of AVC. At any price above this threshold, the firm chooses the quantity where $p=MC$. Market supply is the horizontal sum of individual firms' supply curves.

Supply is not simply a fixed quantity that firms ``have.'' It is a relationship between the price a firm can receive and the profit-maximising quantity it chooses to sell.

\section*{Chapter Review}

\subsection*{Key concepts and equations}
\begin{itemize}
  \item $TC=FC+VC$.
  \item $ATC=AFC+AVC$.
  \item $MC=dTC/dq$.
  \item A competitive firm chooses $p=MC$ when operating.
  \item The short-run shutdown condition for producing is $p\geq \min AVC$.
  \item The short-run supply curve is the MC curve above minimum AVC.
\end{itemize}

\section*{Exercises}

\subsection*{Conceptual questions}
\begin{enumerate}[label=\arabic*.]
  \item Distinguish fixed and variable cost. Give two examples of each for a restaurant.
  \item Why does marginal cost not depend on fixed cost?
  \item Explain why MC crosses AVC and ATC at their minimum points.
  \item Distinguish the short-run shutdown decision from the long-run exit decision.
\end{enumerate}

\subsection*{Graphical exercises}
\begin{enumerate}[resume,label=\arabic*.]
  \item Draw AFC, AVC, ATC, and MC curves with their standard shapes. Mark the minimum points of AVC and ATC.
  \item On the same diagram, show a competitive firm facing three prices: one below minimum AVC, one between minimum AVC and ATC, and one above ATC.
\end{enumerate}

\subsection*{Analytical exercises}
\begin{enumerate}[resume,label=\arabic*.]
  \item A firm's cost function is $C(q)=100+4q+q^2$. Find $FC$, $VC$, $AVC$, $ATC$, and $MC$.
  \item For the firm in Exercise 7, find the output chosen when market price is $20$. Does the firm produce when price is $3$? Explain.
  \item A competitive firm faces $p=15$ and has $TC(q)=50+3q+q^2$. Calculate output, revenue, total cost, and profit.
\end{enumerate}

\subsection*{Exam-style multiple-choice questions}
\begin{enumerate}[resume,label=\arabic*.]
  \item A competitive firm's short-run supply curve is: (a) its demand curve; (b) its MC curve above minimum AVC; (c) its ATC curve; or (d) its fixed-cost curve.
  \item A firm should shut down in the short run if: (a) $p<ATC$; (b) $p<AVC$; (c) $p<MC$ at every output; or (d) fixed cost is positive.
  \item If marginal cost is below average total cost, producing an additional unit will: (a) raise ATC; (b) lower ATC; (c) leave ATC unchanged; or (d) make fixed cost zero.
\end{enumerate}

\subsection*{PPE discussion question}

\noindent A firm may continue operating while making an accounting loss. Discuss why this can be economically rational and how this distinction should affect public debates about subsidies to firms or industries.
